% Options for packages loaded elsewhere
\PassOptionsToPackage{unicode}{hyperref}
\PassOptionsToPackage{hyphens}{url}
\documentclass[
]{article}
\usepackage{xcolor}
\usepackage[margin=1in]{geometry}
\usepackage{amsmath,amssymb}
\setcounter{secnumdepth}{-\maxdimen} % remove section numbering
\usepackage{iftex}
\ifPDFTeX
  \usepackage[T1]{fontenc}
  \usepackage[utf8]{inputenc}
  \usepackage{textcomp} % provide euro and other symbols
\else % if luatex or xetex
  \usepackage{unicode-math} % this also loads fontspec
  \defaultfontfeatures{Scale=MatchLowercase}
  \defaultfontfeatures[\rmfamily]{Ligatures=TeX,Scale=1}
\fi
\usepackage{lmodern}
\ifPDFTeX\else
  % xetex/luatex font selection
\fi
% Use upquote if available, for straight quotes in verbatim environments
\IfFileExists{upquote.sty}{\usepackage{upquote}}{}
\IfFileExists{microtype.sty}{% use microtype if available
  \usepackage[]{microtype}
  \UseMicrotypeSet[protrusion]{basicmath} % disable protrusion for tt fonts
}{}
\makeatletter
\@ifundefined{KOMAClassName}{% if non-KOMA class
  \IfFileExists{parskip.sty}{%
    \usepackage{parskip}
  }{% else
    \setlength{\parindent}{0pt}
    \setlength{\parskip}{6pt plus 2pt minus 1pt}}
}{% if KOMA class
  \KOMAoptions{parskip=half}}
\makeatother
\usepackage{graphicx}
\makeatletter
\newsavebox\pandoc@box
\newcommand*\pandocbounded[1]{% scales image to fit in text height/width
  \sbox\pandoc@box{#1}%
  \Gscale@div\@tempa{\textheight}{\dimexpr\ht\pandoc@box+\dp\pandoc@box\relax}%
  \Gscale@div\@tempb{\linewidth}{\wd\pandoc@box}%
  \ifdim\@tempb\p@<\@tempa\p@\let\@tempa\@tempb\fi% select the smaller of both
  \ifdim\@tempa\p@<\p@\scalebox{\@tempa}{\usebox\pandoc@box}%
  \else\usebox{\pandoc@box}%
  \fi%
}
% Set default figure placement to htbp
\def\fps@figure{htbp}
\makeatother
% definitions for citeproc citations
\NewDocumentCommand\citeproctext{}{}
\NewDocumentCommand\citeproc{mm}{%
  \begingroup\def\citeproctext{#2}\cite{#1}\endgroup}
\makeatletter
 % allow citations to break across lines
 \let\@cite@ofmt\@firstofone
 % avoid brackets around text for \cite:
 \def\@biblabel#1{}
 \def\@cite#1#2{{#1\if@tempswa , #2\fi}}
\makeatother
\newlength{\cslhangindent}
\setlength{\cslhangindent}{1.5em}
\newlength{\csllabelwidth}
\setlength{\csllabelwidth}{3em}
\newenvironment{CSLReferences}[2] % #1 hanging-indent, #2 entry-spacing
 {\begin{list}{}{%
  \setlength{\itemindent}{0pt}
  \setlength{\leftmargin}{0pt}
  \setlength{\parsep}{0pt}
  % turn on hanging indent if param 1 is 1
  \ifodd #1
   \setlength{\leftmargin}{\cslhangindent}
   \setlength{\itemindent}{-1\cslhangindent}
  \fi
  % set entry spacing
  \setlength{\itemsep}{#2\baselineskip}}}
 {\end{list}}
\usepackage{calc}
\newcommand{\CSLBlock}[1]{\hfill\break\parbox[t]{\linewidth}{\strut\ignorespaces#1\strut}}
\newcommand{\CSLLeftMargin}[1]{\parbox[t]{\csllabelwidth}{\strut#1\strut}}
\newcommand{\CSLRightInline}[1]{\parbox[t]{\linewidth - \csllabelwidth}{\strut#1\strut}}
\newcommand{\CSLIndent}[1]{\hspace{\cslhangindent}#1}
\setlength{\emergencystretch}{3em} % prevent overfull lines
\providecommand{\tightlist}{%
  \setlength{\itemsep}{0pt}\setlength{\parskip}{0pt}}
\usepackage{fontspec}
\usepackage{graphicx}
\usepackage{pdfpages}
\usepackage{fancyhdr}
\setmainfont{Arial}
\hyphenpenalty=10000
\exhyphenpenalty=10000
\pagestyle{fancy}
\fancyhf{}
\fancyfoot[R]{Page \thepage}
\renewcommand{\headrulewidth}{0pt}
\renewcommand{\footrulewidth}{0pt}
\usepackage{bookmark}
\IfFileExists{xurl.sty}{\usepackage{xurl}}{} % add URL line breaks if available
\urlstyle{same}
\hypersetup{
  hidelinks,
  pdfcreator={LaTeX via pandoc}}

\author{}
\date{\vspace{-2.5em}}

\begin{document}

\begin{center}
\textit{\Large Supplementary Data for:}
\end{center}

\begin{center}
\textbf{\Large Recipient Plasma Metabolomics as a Predictor of Heart Transplant Severe Primary Graft Dysfunction}
\end{center}

\vspace{1em}

Joshua D. Preston\textsuperscript{1,2}, Clayton J.
Rust\textsuperscript{1}, Aubrey C. Reed\textsuperscript{1}, Supreet S.
Randhawa\textsuperscript{1}, Ailin Tang\textsuperscript{1}, Kimberly A.
Rooney\textsuperscript{3}, Mikhael Kozmava\textsuperscript{3}, Jaclyn
Weinberg\textsuperscript{2}, Lucy Avant\textsuperscript{1}, Catherine
L.B. McGeoch\textsuperscript{1}, Gustavo A. Parrilla\textsuperscript{1},
Jiada Zhan\textsuperscript{2}, ViLinh Tran\textsuperscript{2}, Michael
E. Halkos\textsuperscript{1}, Muath M. Bishawi\textsuperscript{1}, Mani
A. Daneshmand\textsuperscript{1}, Arshed A. Quyyumi\textsuperscript{3},
Charles D. Searles\textsuperscript{3}, Young-Mi Go\textsuperscript{2},
Dean P. Jones\textsuperscript{2}, Joshua L. Chan\textsuperscript{1}

\vspace{1em}

\textsuperscript{1}Division of Cardiothoracic Surgery, Department of
Surgery, Emory University School of Medicine, Atlanta, GA, USA

\textsuperscript{2}Division of Pulmonary, Allergy, Critical Care, and
Sleep Medicine, Department of Medicine, Emory University School of
Medicine, Atlanta, GA, USA

\textsuperscript{3}Division of Cardiology, Department of Medicine, Emory
University School of Medicine, Atlanta, GA, USA

\vfill

\section{Table of Contents}\label{table-of-contents}

\textbf{Supplementary Figure S1} \dotfill 2

\textbf{Supplementary Figure S2.1} \dotfill 3

\textbf{Supplementary Figure S2.2} \dotfill 4

\textbf{Supplementary Figure S2.3} \dotfill 5

\textbf{Supplementary Figure S2.4} \dotfill 6

\textbf{Supplementary Figure S2.5} \dotfill 7

\textbf{Expanded Metabolomics Methods and Other Technical Details}
\dotfill 8-9

\quad \textit{Sample Collection and Processing} \dotfill 8

\quad \textit{Metabolomics Profiling with Liquid Chromatography--Mass Spectrometry}
\dotfill 8

\quad \textit{Feature Extraction, Annotation, and Identification}
\dotfill 8-9

\quad \textit{Quality Control and Data Transformation} \dotfill 9

\textbf{Statistical Analysis and Data Visualization} \dotfill 9-10

\quad \textit{Partial Least Squares Discriminant Analysis and Volcano Plots}
\dotfill 9

\quad \textit{Pathway Enrichment Analysis} \dotfill 9-10

\quad \textit{Identified and Annotated Metabolite Analysis} \dotfill 10

\quad \textit{Sensitivity Analysis} \dotfill 10

\textbf{References} \dotfill 11-12

\newpage

\includepdf[pages=-,pagecommand={\noindent\textbf{\Large Supplementary Figure S1}\vspace{0.5cm}}]{Components/Figures/PDF/S1.pdf}

\includepdf[pages=-,pagecommand={\noindent\textbf{\Large Supplementary Figure S2.1}\vspace{0.2cm}\\\noindent\textit{\fontsize{11}{13}\selectfont Markers on x-axis indicate Y/N severe PGD status.}\vspace{0.3cm}}]{Components/Figures/PDF/S2.1.pdf}

\includepdf[pages=-,pagecommand={\noindent\textbf{\Large Supplementary Figure S2.2}\vspace{0.2cm}\\\noindent\textit{\fontsize{11}{13}\selectfont Markers on x-axis indicate Y/N severe PGD status.}\vspace{0.3cm}}]{Components/Figures/PDF/S2.2.pdf}

\includepdf[pages=-,pagecommand={\noindent\textbf{\Large Supplementary Figure S2.3}\vspace{0.2cm}\\\noindent\textit{\fontsize{11}{13}\selectfont Markers on x-axis indicate Y/N severe PGD status.}\vspace{0.3cm}}]{Components/Figures/PDF/S2.3.pdf}

\includepdf[pages=-,pagecommand={\noindent\textbf{\Large Supplementary Figure S2.4}\vspace{0.2cm}\\\noindent\textit{\fontsize{11}{13}\selectfont Markers on x-axis indicate Y/N severe PGD status.}\vspace{0.3cm}}]{Components/Figures/PDF/S2.4.pdf}

\includepdf[pages=-,pagecommand={\noindent\textbf{\Large Supplementary Figure S2.5}\vspace{0.5cm}}]{Components/Figures/PDF/S2.5.pdf}

\section{\texorpdfstring{\textbf{Expanded Metabolomics Methods and Other
Technical
Details}}{Expanded Metabolomics Methods and Other Technical Details}}\label{expanded-metabolomics-methods-and-other-technical-details}

In keeping with consensus reporting standards for metabolomics
experiments,\textsuperscript{1} we provide in-depth details as follows:

\subsection{\texorpdfstring{\emph{Sample Collection and
Processing}}{Sample Collection and Processing}}\label{sample-collection-and-processing}

Blood samples were collected immediately before transplant into K2EDTA
plasma preparation tubes. Once samples were collected from patients,
they were immediately placed on ice. The samples remained on ice and
were centrifuged at 2500 x g for 15 minutes at 4°C within 5 hours of
collection (the mean time from collection to processing for all samples
= 1 hour). Following centrifugation, three aliquots were made from the
supernatant, which were stored at -80°C shortly after.

\subsection{\texorpdfstring{\emph{Metabolomics Profiling with Liquid
Chromatography--Mass
Spectrometry}}{Metabolomics Profiling with Liquid Chromatography--Mass Spectrometry}}\label{metabolomics-profiling-with-liquid-chromatographymass-spectrometry}

High-resolution metabolomics was performed in the Emory Clinical
Biomarkers Laboratory as previously described.\textsuperscript{2} Plasma
samples were thawed, and 10 μL of plasma was added to 190 μL of
acetonitrile containing 2\% internal standards. Samples were then
vortexed and allowed to incubate on ice for 30 minutes. Following
incubation, samples were centrifuged at 20,817 x g for 10 minutes at
4°C. Supernatants were then transferred into autosampler vials. These
were loaded onto autosampler plates and were kept at -20°C until
analysis. Pooled control plasma (Equitech-Bio, SHP45) and National
Institute of Standards and Technology (NIST) reference plasma were
prepared and analyzed in tandem.

\noindent\hspace{1.5em}Metabolomics analysis was conducted with a liquid
chromatography--mass spectrometry system (LC--MS) consisting of a
Vanquish Duo UHPLC coupled to an Orbitrap ID-X Tribrid Mass Spectrometer
(Thermo Scientific). Autosampler plates containing samples were
maintained at 4°C in the autosampler of the LC--MS throughout the
analysis. Samples were analyzed in triplicate on a 5-min method using
C18 chromatography coupled to negative electrospray ionization (ESI)
(C18-) and hydrophilic interaction liquid chromatography coupled to
positive ESI (HILIC+). Analyte separation for HILIC was performed with a
Waters Acquity BEH Amide HILIC column (2.1 mm x 100 mm, 1.7 μm particle
size) and gradient elution with LC--MS-grade solvents and additives. The
HILIC mobile phases included (Buffer A) water with 1 mM ammonium acetate
and 0.1\% formic acid and (Buffer B) 95\% acetonitrile with 1 mM
ammonium acetate and 0.1\% formic acid. For the HILIC gradient, an
initial 0.5 min hold at 90\% B was followed by a linear decrease to 20\%
B from 0.6 to 2.55 min, a 2 min hold, and a 5-min re-equilibration
period. C18 chromatography was performed on a Thermo Hypersil Gold C18
column (2.1 mm x 100 mm, 1.9 μm particle size). The C18 mobile phases
included (Buffer A) water with 1 mM ammonium acetate and (Buffer B) 99\%
acetonitrile with 1 mM ammonium acetate. For the C18 gradient, an
initial 0.5 min hold at 1\% B was followed by a 0.75 min linear increase
to 99\% B, held for 3.75 min, and a 5-min re-equilibration period. The
flow rate for both methods was 0.3 mL/min, and the column compartment
was heated to 45 °C. The mass spectrometer was operated at 120k
resolution and MS1 scans were collected for \textit{m/z} 85-1,275. Tune
parameters consisted of sheath gas at 50, auxiliary gas at 10, and sweep
gas at 1. The spray voltage was set to 3.50 kV for ESI+ and -2.75 kV for
ESI-.

\subsection{\texorpdfstring{\emph{Feature Extraction, Annotation, and
Identification}}{Feature Extraction, Annotation, and Identification}}\label{feature-extraction-annotation-and-identification}

ProteoWizard v3\textsuperscript{3} was used to convert raw spectral
files to .mzXML files. Untargeted feature tables were then generated
from extracted data using established laboratory workflows,
software\textsuperscript{4--6} (namely, apLCMS and xMSAnalyzer), and an
in-house R package, \textit{MetaboJanitoR}, which is available upon
request. All feature tables were corrected for batch effects using
ComBat\textsuperscript{7} from the \textit{sva}
package.\textsuperscript{8}

\noindent\hspace{1.5em}Untargeted feature tables were comprised of all
HILIC+ and C18- features observed with a unique mass-to-charge ratio
(\textit{m/z}) and corresponding retention time (RT).
Identified/annotated feature tables were made up of Level
1\textsuperscript{9} and Level 3\textsuperscript{9} identifications.
Level 1 identifications were obtained by ion dissociation mass
spectrometry relative to authentic standards (mass error = ± 5 ppm,
retention time threshold = ± 30 seconds). Level 3 identifications were
made using an in-house annotation algorithm (see
\url{https://github.com/jamesjiadazhan} for the future release of this R
package upon its publication), which generates annotations based on
precursor-product correlations (note: this version of the algorithm
chose a single best adduct for the feature table based on the strength
of the precursor-product correlation alone). All raw data files (.mzXML
format) and untargeted feature tables analyzed for this study will be
available on October 31, 2026, through the NIH Common Fund's National
Metabolomics Data Repository (NMDR) website, Metabolomics Workbench
(Project ID PR002742; Study ID ST004328; Project DOI
\url{http://dx.doi.org/10.21228/M87549}).

\subsection{\texorpdfstring{\emph{Quality Control and Data
Transformation}}{Quality Control and Data Transformation}}\label{quality-control-and-data-transformation}

Features that were not detected in at least 80\% of samples were
excluded from analysis, but were included as background for pathway
enrichment analysis. All `0' values (indicating non-detection by the
instrument) were replaced with ½ the minimum value for that feature
detected across all samples.\textsuperscript{10} Following this, all
intensity (peak area) values were log\(_2\)-transformed prior to further
analysis. Features in identified/annotated feature tables that
significantly differed between study groups (Severe PGD versus No Severe
PGD) were manually inspected and evaluated for biological plausibility
and annotation quality control. Only annotated features encompassing
endogenous, microbiome-derived, or relevant pharmacological metabolites
were retained for analysis and visualization.

\section{\texorpdfstring{\textbf{Statistical Analysis and Data
Visualization}}{Statistical Analysis and Data Visualization}}\label{statistical-analysis-and-data-visualization}

The complete analytical pipeline and source code used for data
processing and statistical analyses are publicly available at
\url{https://github.com/jdpreston30/PGD-preop-metabolomics}. A
containerized Docker image of the complete computational environment is
available at Docker Hub
(\url{https://hub.docker.com/r/jdpreston30/pgd-preop-metabolomics}), and
the archived analysis pipeline is permanently available at Zenodo
(\url{https://doi.org/10.5281/zenodo.17875844}).

\noindent\hspace{1.5em}Data were analyzed and visualized using Microsoft
Excel (Mac v16.101.294, Microsoft Corporation, 2025), R
(v4.5.1),\textsuperscript{11} and MetaboAnalyst
(v6.0)\textsuperscript{12} via the \textit{MetaboAnalystR}
package.\textsuperscript{13,14} All code was written, compiled, and run
in Visual Studio Code (v1.104.2, Microsoft Corporation, 2025). All data
visualizations were created using \textit{ggplot2},\textsuperscript{15}
\textit{igraph},\textsuperscript{16--18} and
\textit{ggraph}.\textsuperscript{19} All figures were compiled using
\textit{cowplot}.\textsuperscript{20} Epidemiologic statistics were
performed using the \textit{TernTablesR} package.\textsuperscript{21}
For all analyses, statistical significance was set at α = .05.

\subsection{\texorpdfstring{\emph{Partial Least Squares Discriminant
Analysis and Volcano
Plots}}{Partial Least Squares Discriminant Analysis and Volcano Plots}}\label{partial-least-squares-discriminant-analysis-and-volcano-plots}

Partial Least Squares Discriminant Analysis (PLS-DA) was performed using
the R package, \textit{mixOmics}.\textsuperscript{22} Volcano plots were
created by calculating \textit{P} values (based on Welch's
\textit{t}-tests on log\(_2\)-transformed data) to compare the means of
each (untargeted) feature between groups, while fold changes were
determined using untransformed data. All \textit{P} value calculations
were performed using R's \textit{stats} package. For purposes of
visualization, \textit{P} values were then -log\(_{10}\)-transformed,
and fold changes were log\(_2\)-transformed.

\subsection{\texorpdfstring{\emph{Pathway Enrichment
Analysis}}{Pathway Enrichment Analysis}}\label{pathway-enrichment-analysis}

All pathway enrichment analysis was performed using
Mummichog.\textsuperscript{23} A full list of \textit{P} values
comparing means of all features between groups was first generated in R
using the \textit{stats} package via Welch's \textit{t}-tests. These
features, including the \textit{m/z}, retention time, and ESI mode,
along with their corresponding \textit{P} values, were then analyzed
using the `Functional Analysis' tool for LC--MS from
MetaboAnalyst\textsuperscript{12} (run via the \textit{MetaboAnalystR}
package\textsuperscript{13}), which utilizes
Mummichog\textsuperscript{23} as its core algorithm. Notably, this tool
allows for the usage of Mummichog v2.0, which takes into account
retention time and ESI mode, whereas the original algorithm simply used
\textit{m/z}. The parameters for the pathway enrichment included the
following: mixed mode, 5 ppm mass tolerance, primary ions enforced,
\textit{P} value cutoff = default for top 10\% of peaks, KEGG and MFN
\textit{homo sapiens} databases. A bubble plot was then constructed
based on the \textit{P} values (Fisher's) and enrichment factors
generated by the algorithm output.

\subsection{\texorpdfstring{\emph{Identified and Annotated Metabolite
Analysis}}{Identified and Annotated Metabolite Analysis}}\label{identified-and-annotated-metabolite-analysis}

For identified/annotated metabolites, Welch's \textit{t}-tests were
employed to determine the \textit{P} values comparing means between each
group. FDR adjustments were performed using the Benjamini-Hochberg (BH)
procedure. A representative set of various families of metabolites was
selected and displayed in the main text figures (\textbf{​Fig. 3​A-D}),
but all significantly differing annotated or identified features are
displayed in \textbf{Supplementary Fig. S2}. Both raw \textit{P} values
and FDR-adjusted \textit{q} values are displayed on all plots displaying
annotated/identified metabolites.

\subsection{\texorpdfstring{\emph{Sensitivity
Analysis}}{Sensitivity Analysis}}\label{sensitivity-analysis}

Given the exploratory and hypothesis-generating nature of the untargeted
metabolomics study performed herein, which measured over 20,000
individual metabolic features, a conventional power analysis is not
applicable. However, we conducted a quantitative post-hoc sensitivity
analysis (further details of which can be found in the GitHub
repository) to determine the minimum detectable effect size based on the
group sizes in our study. With 62 patients (8 with severe PGD and 54
without), we had \textasciitilde80\% power (α = .05) to detect a minimum
standardized mean difference of Cohen's d ≈ 1.08. Based on observed
within-group variance (calculated using the filtered, i.e., post-QC,
untargeted feature table), the median minimum detectable fold-change was
1.87 (IQR 1.46--2.99) across all features. Considering this, this study
was powered to detect very large metabolic differences between groups,
which is consistent with the exploratory and hypothesis-generating
design. More subtle differences would require larger cohorts in future
studies.

\newpage

\section*{\texorpdfstring{\textbf{References}}{References}}\label{references}
\addcontentsline{toc}{section}{\textbf{References}}

\protect\phantomsection\label{refs}
\begin{CSLReferences}{0}{1}
\bibitem[\citeproctext]{ref-sumnerProposedMinimumReporting2007}
\CSLLeftMargin{1. }%
\CSLRightInline{Sumner L.W., Amberg A., Barrett D., et al. Proposed
minimum reporting standards for chemical analysis: {Chemical Analysis
Working Group} ({CAWG}) {Metabolomics Standards Initiative} ({MSI}).
\emph{Metabolomics}. 2007;3(3):211-221.
doi:\href{https://doi.org/10.1007/s11306-007-0082-2}{10.1007/s11306-007-0082-2}}

\bibitem[\citeproctext]{ref-liuReferenceStandardizationQuantification2020}
\CSLLeftMargin{2. }%
\CSLRightInline{Liu K.H., Nellis M., Uppal K., et al. Reference
standardization for quantification and harmonization of large-scale
metabolomics. \emph{Analytical Chemistry}. 2020;92(13):8836-8844.
doi:\href{https://doi.org/10.1021/acs.analchem.0c00338}{10.1021/acs.analchem.0c00338}}

\bibitem[\citeproctext]{ref-chambersCrossplatformToolkitMass2012}
\CSLLeftMargin{3. }%
\CSLRightInline{Chambers M.C., Maclean B., Burke R., et al. A
cross-platform toolkit for mass spectrometry and proteomics.
\emph{Nature Biotechnology}. 2012;30(10):918-920.
doi:\href{https://doi.org/10.1038/nbt.2377}{10.1038/nbt.2377}}

\bibitem[\citeproctext]{ref-uppal2013}
\CSLLeftMargin{4. }%
\CSLRightInline{Uppal K., Soltow Q.A., Strobel F.H., et al.
{xMSanalyzer}: {automated} pipeline for improved feature detection and
downstream analysis of large-scale, non-targeted metabolomics data.
\emph{BMC Bioinformatics}. 2013;14(1):15.
doi:\href{https://doi.org/10.1186/1471-2105-14-15}{10.1186/1471-2105-14-15}}

\bibitem[\citeproctext]{ref-jarrellPlasmaAcylcarnitineLevels2020}
\CSLLeftMargin{5. }%
\CSLRightInline{Jarrell Z.R., Smith M.R., Hu X., et al. Plasma
acylcarnitine levels increase with healthy aging. \emph{Aging}.
2020;12(13):13555-13570.
doi:\href{https://doi.org/10.18632/aging.103462}{10.18632/aging.103462}}

\bibitem[\citeproctext]{ref-yuApLCMSAdaptiveProcessing2009}
\CSLLeftMargin{6. }%
\CSLRightInline{Yu T., Park Y., Johnson J.M., Jones D.P.
{apLCMS}---{adaptive} processing of high-resolution {LC/MS} data.
\emph{Bioinformatics}. 2009;25(15):1930-1936.
doi:\href{https://doi.org/10.1093/bioinformatics/btp291}{10.1093/bioinformatics/btp291}}

\bibitem[\citeproctext]{ref-johnson2007}
\CSLLeftMargin{7. }%
\CSLRightInline{Johnson W.E., Li C., Rabinovic A. Adjusting batch
effects in microarray expression data using empirical {Bayes} methods.
\emph{Biostatistics (Oxford, England)}. 2007;8(1):118-127.
doi:\href{https://doi.org/10.1093/biostatistics/kxj037}{10.1093/biostatistics/kxj037}}

\bibitem[\citeproctext]{ref-leek2012}
\CSLLeftMargin{8. }%
\CSLRightInline{Leek J.T., Johnson W.E., Parker H.S., Jaffe A.E., Storey
J.D. The sva package for removing batch effects and other unwanted
variation in high-throughput experiments. \emph{Bioinformatics (Oxford,
England)}. 2012;28(6):882-883.
doi:\href{https://doi.org/10.1093/bioinformatics/bts034}{10.1093/bioinformatics/bts034}}

\bibitem[\citeproctext]{ref-schymanskiIdentifyingSmallMolecules2014}
\CSLLeftMargin{9. }%
\CSLRightInline{Schymanski E.L., Jeon J., Gulde R., et al. Identifying
small molecules via high resolution mass spectrometry: Communicating
confidence. \emph{Environmental Science \& Technology}.
2014;48(4):2097-2098.
doi:\href{https://doi.org/10.1021/es5002105}{10.1021/es5002105}}

\bibitem[\citeproctext]{ref-weiMissingValueImputation2018}
\CSLLeftMargin{10. }%
\CSLRightInline{Wei R., Wang J., Su M., et al. Missing value imputation
approach for mass spectrometry-based metabolomics data. \emph{Scientific
Reports}. 2018;8(1):663.
doi:\href{https://doi.org/10.1038/s41598-017-19120-0}{10.1038/s41598-017-19120-0}}

\bibitem[\citeproctext]{ref-rcoreteam2025}
\CSLLeftMargin{11. }%
\CSLRightInline{{R}: A language and environment for statistical
computing. {R Core Team}; 2025. \url{https://www.R-project.org/}.
{[}Accessed 10 December 2025{]}}

\bibitem[\citeproctext]{ref-pangMetaboAnalyst60Unified2024}
\CSLLeftMargin{12. }%
\CSLRightInline{Pang Z., Lu Y., Zhou G., et al. {MetaboAnalyst} 6.0:
{towards} a unified platform for metabolomics data processing, analysis
and interpretation. \emph{Nucleic Acids Research}.
2024;52(W1):W398-W406.
doi:\href{https://doi.org/10.1093/nar/gkae253}{10.1093/nar/gkae253}}

\bibitem[\citeproctext]{ref-xia2025}
\CSLLeftMargin{13. }%
\CSLRightInline{{MetaboAnalystR} 4.2: {a} unified {LC-MS/MS} workflow
for global metabolomics and exposomics. {Xia Lab}; 2025.
\url{https://github.com/xia-lab/MetaboAnalystR}. {[}Accessed 10 December
2025{]}}

\bibitem[\citeproctext]{ref-MetaboAnalystRv3}
\CSLLeftMargin{14. }%
\CSLRightInline{Pang Z., Chong J., Li S., Xia J. {MetaboAnalystR} 3.0:
Toward an optimized workflow for global metabolomics.
\emph{Metabolites}. 2020.
doi:\href{https://doi.org/doi:\%2010.3390/metabo10050186}{doi:
10.3390/metabo10050186}}

\bibitem[\citeproctext]{ref-wickhamGgplot2ElegantGraphics2016}
\CSLLeftMargin{15. }%
\CSLRightInline{Wickham H. \emph{ggplot2: Elegant graphics for data
analysis}. Second edition. Springer; 2016.}

\bibitem[\citeproctext]{ref-csardiIgraphSoftwarePackage2006}
\CSLLeftMargin{16. }%
\CSLRightInline{Csárdi G., Nepusz T. The igraph software package for
complex network research. \emph{InterJournal}. 2006;Complex
Systems:1695. \url{https://igraph.org}}

\bibitem[\citeproctext]{ref-antonovIgraphEnablesFast2023}
\CSLLeftMargin{17. }%
\CSLRightInline{Antonov M., Csárdi G., Horvát S., et al. {igraph}
enables fast and robust network analysis across programming languages.
\emph{arXiv preprint arXiv:231110260}. 2023.
doi:\href{https://doi.org/10.48550/arXiv.2311.10260}{10.48550/arXiv.2311.10260}}

\bibitem[\citeproctext]{ref-csardiIgraphNetworkAnalysis2025}
\CSLLeftMargin{18. }%
\CSLRightInline{Csárdi G., Nepusz T., Traag V., et al. {igraph}: Network
analysis and visualization in {R}. 2025.
doi:\href{https://doi.org/10.5281/zenodo.7682609}{10.5281/zenodo.7682609}}

\bibitem[\citeproctext]{ref-pedersenGgraphImplementationGrammar2025}
\CSLLeftMargin{19. }%
\CSLRightInline{Pedersen T.L. {ggraph}: An implementation of grammar of
graphics for graphs and networks. 2025.
\url{https://ggraph.data-imaginist.com}}

\bibitem[\citeproctext]{ref-clauso.wilke2025}
\CSLLeftMargin{20. }%
\CSLRightInline{Wilke C.O. {cowplot}: Streamlined plot theme and plot
annotations for {'}ggplot2{'}. 2025.}

\bibitem[\citeproctext]{ref-preston2025}
\CSLLeftMargin{21. }%
\CSLRightInline{Preston J.D. {TernTablesR}: Automated statistics and
table generation for clinical research in {R}. 2025.
\url{https://github.com/jdpreston30/TernTablesR}. {[}Accessed 10
December 2025{]}}

\bibitem[\citeproctext]{ref-rohart2017}
\CSLLeftMargin{22. }%
\CSLRightInline{Rohart F., Gautier B., Singh A., Lê Cao K.A. {mixOmics}:
An {R} package for `omics feature selection and multiple data
integration. Schneidman D., ed. \emph{PLOS Computational Biology}.
2017;13(11):e1005752.
doi:\href{https://doi.org/10.1371/journal.pcbi.1005752}{10.1371/journal.pcbi.1005752}}

\bibitem[\citeproctext]{ref-li2013}
\CSLLeftMargin{23. }%
\CSLRightInline{Li S., Park Y., Duraisingham S., et al. Predicting
network activity from high throughput metabolomics. \emph{PLoS
computational biology}. 2013;9(7):e1003123.
doi:\href{https://doi.org/10.1371/journal.pcbi.1003123}{10.1371/journal.pcbi.1003123}}

\end{CSLReferences}

\end{document}
